\documentclass[12pt]{article}

\usepackage{sbc}
\usepackage{graphicx,url}
\usepackage[utf8]{inputenc}
\usepackage[english]{babel}
\usepackage{amsmath}

\sloppy

\title{Predicting mutational effects in proteins using pre-trained\\Protein Language Model embeddings}

\author{Douglas A. C. Canevarollo\inst{1}}

\address{
  Laboratório de Bioinformática\\
  Universidade Estadual Paulista (UNESP)\\Instituto de Biociências, Letras e Ciências Exatas (IBILCE)\\São José do Rio Preto, SP -- Brazil
  \email{douglas.canevarollo@unesp.br}
}

\begin{document}

\maketitle

\begin{abstract}
Single-nucleotide variants in the human exome, especially missense mutations, can severely disrupt protein structure and function and are linked to rare diseases, complex disorders, and cancer. Predicting their pathogenicity is challenging due to the large number of variants per individual and limited experimental data, particularly for genes associated with rare diseases. ClinVar provides curated clinical labels, while gnomAD supplies population-scale allele frequencies that help distinguish rare, potentially pathogenic variants from common, likely benign ones. Protein Language Models (pLMs), such as ESM-2 and Prot-T5, have recently been used to predict the functional impact of human missense variants without task-specific supervision. This paper evaluates and compares two zero-shot mutation representation strategies --- log-ratio scores derived from model output probabilities, and delta-embedding norms derived from hidden-state shifts --- across both pLMs, for predicting pathogenicity-linked functional impact in a per-gene low-data regime. An explicit per-gene evaluability criterion is used to define this regime, and all four representation--model combinations are assessed using ROC-AUC and PR-AUC. Results show that the ESM-2 log-ratio representation substantially outperforms all other combinations (median ROC-AUC of 0.951), while delta-embedding norms perform considerably worse and, for ESM-2 specifically, fall below the level expected by chance --- a finding discussed as an informative architectural contrast rather than an experimental failure. The current analysis is restricted to variants on chromosome 1 of gnomAD; extension to additional chromosomes and the incorporation of population allele frequencies as a model feature are left as directions for future work.
\end{abstract}

\section{Introduction}

Single-nucleotide variants in the human exome, especially missense mutations, can drastically alter the structure and function of the target protein, and are often associated with cancer, complex diseases, and rare genetic disorders. Hence, the accurate classification of these variants as pathogenic or benign stands as a central challenge in genomic medicine. On average, the human exome harbors around 20,000 single-nucleotide variants \cite{alirezaie}, further amplifying the need for scalable computational tools to support clinical databases. Errors in variant classification may directly impact medical diagnosis and treatment recommendations \cite{sharo}.

In this context, population genomics databases have become key resources for pathogenicity prediction \cite{sharo}. The Genome Aggregation Database (gnomAD) \cite{gnomad}, for instance, aggregates variant frequencies across large cohorts of healthy individuals, enabling the identification of rare variants --- often potentially pathogenic --- while filtering out common variants unlikely to cause rare diseases. Such databases can be combined with specialized annotation resources, as ClinVar \cite{clinvar}, which curates clinical interpretations of variants. Pathogenicity prediction tools such as ClinPred \cite{alirezaie} already leverage this combination of ClinVar and gnomAD to train and validate their classification models.

On another front, advances in Artificial Intelligence (AI) have enabled Large Language Models (LLMs) to revolutionize Natural Language Processing (NLP). Their adaptation to protein sequences has led to the development of powerful protein Language Models (pLMs) \cite{prottrans}. These models have been used to predict the impact of missense mutations in human proteins, including architectures such as Evolutionary Scale Modeling (ESM-2) \cite{esm2} and Prot-T5 \cite{prottrans}, in tasks ranging from protein stability to ``variant effect'' prediction related to phenotypes.

However, although these pLMs are promising, few systematic studies evaluate specific representation strategies under scenarios of low data availability per gene, using ClinVar as a label source and gnomAD as a population filter \cite{stein}. Representation choices --- such as pLM log-ratio scores versus delta-embedding distances --- may have a larger impact on performance than the choice of the pLM itself \cite{alirezaie}, and the relative contribution of population-level features such as gnomAD allele frequency to this comparison remains comparatively underexplored. Moreover, for many genes associated with rare diseases, the number of experimentally or clinically annotated variants is small, which reinforces the need to study zero-shot and low-data settings \cite{sharo}.

In this light, this paper aims to evaluate and compare zero-shot representation strategies based on protein Language Models for predicting the functional impact linked to the pathogenicity of human missense variants. ClinVar is used as the source of pathogenicity labels, and gnomAD is used both to assign population allele frequencies to each variant and to define the chromosome-restricted scope of the current analysis (chromosome 1), focusing especially on low-data scenarios defined at the gene level. Two pLMs --- ESM-2 and Prot-T5 --- and two mutation representations --- log-ratio scores derived from model output probabilities, and delta-embedding distances derived from hidden-state shifts --- are compared under an explicit per-gene evaluability criterion, with no supervised fine-tuning or downstream classifiers involved at any stage.

The main empirical finding is that the ESM-2 log-ratio representation substantially outperforms all other representation--model combinations evaluated, while delta-embedding-based scores perform considerably worse and, in the case of ESM-2, fall below chance-level discrimination --- a result that is examined as an informative contrast between architectures rather than treated as a failure of the method. The incorporation of gnomAD allele frequencies as a predictive feature, alongside expansion beyond chromosome 1, is left as a direction for future work.

\section{Related Work}

The prediction of missense variant effect has long been a central problem in computational biology and genomic medicine. Early approaches relied primarily on evolutionary conservation, multiple sequence alignments, and hand-engineered sequence features to estimate whether an amino acid substitution would be deleterious. Methods such as SIFT, PolyPhen-2, MutPred2, EVmutation, DeepSequence and EVE established that evolutionary information can be highly informative for variant interpretation, but they also revealed an important limitation: their performance often depends on the availability of sufficient homologous sequences or aligned family data. This becomes particularly problematic in genes with sparse evolutionary coverage or limited experimental annotation, which is a common setting in clinical genetics \cite{liu}.

Clinical and population databases have played a crucial role in this research area. ClinVar provides curated interpretations of human variants and is widely used as a source of pathogenicity labels, while gnomAD offers population-scale allele frequency information that helps distinguish rare potentially disease-causing variants from common benign ones. Combining clinical labels with population frequencies has proven useful for benchmarking pathogenicity predictors, but recent work has cautioned that benchmarking and training directly on curated variant sets such as ClinVar can introduce ascertainment bias and data circularity; orthogonal, population-genetics-based benchmarking approaches have been proposed to mitigate this concern \cite{mikhail}. The reliability of these resources still depends on the gene, the disease context, and the distribution of available annotations, which motivates careful evaluation protocols and gene-specific analyses \cite{clinvar}.

The emergence of protein language models has substantially changed the landscape of variant effect prediction. Inspired by large language models in natural language processing, protein language models learn statistical representations from massive collections of amino acid sequences without requiring explicit supervision. These embeddings capture evolutionary and functional regularities that can be transferred to downstream tasks such as structure prediction, function annotation, stability estimation, and mutation effect prediction. Representative models such as ESM-2 \cite{esm2} and Prot-T5 \cite{prottrans} have shown that sequence-only pretraining can produce biologically meaningful representations, even in settings where labeled data are scarce. Foundational work demonstrated that, using only zero-shot inference and without any supervision from experimental data or additional training, protein language models can capture the functional effects of mutations at state-of-the-art levels \cite{meier}, which has made these models especially attractive for prediction problems where conventional supervised methods struggle due to limited labeled data.

A growing body of work has specifically examined how protein language models can be used to predict mutation pathogenicity. One important line of research showed that large pre-trained protein language models act as zero-shot predictors for disease-causing mutations and even clinical prognosis in cancer-related genes, outperforming strong baselines such as EVE in certain benchmarks \cite{liu}. These findings suggest that mutation scores derived from the model's output distribution, such as evolutionary index or log-ratio style measures, may capture functionally relevant constraints directly from sequence statistics. In contrast to purely zero-shot scoring, other approaches have fine-tuned protein language models directly on amino-acid-level functional annotations in order to better separate pathogenic from benign missense variants and to support the reclassification of variants of uncertain significance \cite{saadat}. Other studies have explored alternative mutation representations, including embedding differences between wild-type and mutant sequences, raising the possibility that representation choice may be as influential as model architecture itself for downstream performance --- a question this study addresses directly by comparing both representation families under a common evaluation protocol.

At the same time, recent literature has highlighted the limits of sequence-only modeling. Variant effects are not determined by sequence alone, since structural context, residue interactions, and local protein environments can strongly influence functional consequences. In response, structure-informed protein language models have been proposed to inject structural supervision during training while still using sequence-only inputs at prediction time. A recent study reported that structure-informed models were robust predictors of variant effects across deep mutational scanning benchmarks and that adding structural context could improve robustness more effectively than simply scaling model size or training data \cite{sun}. These results suggest that biological context beyond raw sequence can help reduce overfitting and improve generalization, especially in difficult protein families.

Recent benchmark studies have also stressed the importance of evaluation design. Rather than relying only on global averages, modern assessments increasingly examine calibration, zero-shot behavior, gene-specific performance, and performance under low-data regimes. This is particularly relevant for missense variant prediction, where the number of labeled examples per gene is often small and the class balance can vary substantially \cite{vasquez}. Benchmarking work has shown that continuous and binary predictors may behave differently depending on the target setting, and that interpretation frameworks should consider both predictive quality and practical clinical usability. In this context, low-data experiments are especially valuable because they approximate the real constraints faced in rare-disease and gene-specific applications.

Taken together, the literature indicates a clear progression from conservation-based methods to protein language model approaches, and more recently to structure-informed and benchmark-driven frameworks \cite{sun}. However, there is still room for systematic comparison of mutation representation strategies under low-data conditions, especially when ClinVar labels and gnomAD frequencies are combined at the gene level. This study is positioned within this gap by evaluating different protein language model representations for missense variant pathogenicity prediction in sparse-data settings, with explicit attention to zero-shot regimes.

\section{Materials and Methods}

\subsection{Data Sources and Variant Collection}

Human missense variants were collected from ClinVar and gnomAD. ClinVar served as the primary source of clinical pathogenicity annotations, whereas gnomAD provided population allele frequency information for each variant. The current study is restricted in scope to variants located on chromosome 1, using gnomAD exome data for this chromosome; extension to the remaining chromosomes is left for future work.

ClinVar Variant Call Format (VCF) files were downloaded and functionally annotated with snpEff \cite{snpEff} to obtain transcript-level and protein-level consequence annotations, including amino acid substitutions and protein coordinates. Only variants annotated as missense were retained for downstream analyses.

For each retained variant, the following information was extracted: chromosomal position, reference and alternate alleles, protein position, wild-type amino acid, mutant amino acid, ClinVar clinical significance label, and gnomAD allele frequency. Genomic coordinates were annotated against the GRCh38 \cite{grch38} human reference build, and protein-level consequence annotations were derived from the \textit{ANN} field generated during functional annotation with snpEff.

To reduce label noise, ClinVar annotations were grouped into binary classes: pathogenic and likely pathogenic were assigned to the pathogenic class, whereas benign and likely benign were assigned to the benign class. Variants labeled as uncertain significance, conflicting interpretations, or lacking clear review criteria were excluded from the primary analyses.

After missense filtering and binary label assignment, the resulting dataset comprised 76,990 variants across 8,297 unique genes.

\subsection{Sequence Retrieval and Quality Control}

Full-length protein reference sequences were retrieved from the Ensembl REST API \cite{ensembl} using the transcript identifier annotated for each variant (\texttt{ENST} format) and matched to variant records using the protein-level coordinates obtained during functional annotation. Sequence retrieval succeeded for 97.12\% of annotated transcripts; variants for which no protein sequence could be retrieved were discarded.

Two additional quality-control steps were applied before downstream representation extraction. First, variants whose annotated protein position exceeded the length of the retrieved reference sequence were removed (132 variants, 0.17\% of the dataset). Second, the wild-type amino acid reported by ClinVar/snpEff was cross-checked against the residue present in the retrieved Ensembl sequence at the annotated one-based position; variants for which these did not match were discarded (998 variants, approximately 1.28\% of the dataset), since such a mismatch indicates an unresolved discrepancy between the genomic annotation and the reference transcript used for sequence retrieval rather than a genuine variant.

Sequences containing the character ``X'' (unknown or non-standard amino acid) were not filtered at this stage; this filtering was applied downstream, immediately before protein language model inference (see Section~\ref{sec:pseq}).

\subsection{Population Frequency Integration}

Population allele frequencies were retrieved from gnomAD exome chromosome 1 data and matched to ClinVar variants using genomic coordinates and allelic information; variants without a corresponding gnomAD entry were assigned an allele frequency of zero. Because allele frequency distributions are highly skewed, frequencies were transformed using logarithmic scaling, as shown below:

\begin{equation}
    log\_AF = log_{10}(AF + \epsilon)
\end{equation}

where $\epsilon = 10^{-8}$ is a small constant added to avoid undefined values for zero-frequency variants. The transformed allele frequency ($log\_AF$) was computed and stored for every variant but was not used as a predictive feature in the experiments reported in this paper. Given known concerns about ascertainment bias in ClinVar with respect to allele rarity --- rare variants are disproportionately curated as pathogenic partly because they are rare --- incorporating gnomAD frequency as a model feature is left as a direction for future work.

\subsection{Protein Sequence Representation}
\label{sec:pseq}

Protein representations were generated using two pre-trained protein language models: ESM-2 \cite{esm2}, specifically the \texttt{esm2\_t33\_650M\_UR50D} checkpoint (650M parameters), and Prot-T5 \cite{prottrans}, specifically \texttt{Rostlab/prot\_t5\_xl\_uniref50} (full precision, 3B parameters). For each protein, the wild-type amino acid sequence was provided as input to the corresponding model, and mutation-aware representations were then derived using the strategies described below.

\subsection{Handling of Long Protein Sequences}

A substantial fraction of retrieved protein sequences exceeded the context-length limits practical for the models used: 42.31\% of sequences were longer than 1,024 residues, and 20.74\% of evaluable variants had a protein position beyond residue 1,024. To accommodate this while preserving local sequence context around each mutation, a windowing strategy centered on the mutated position was adopted: for ESM-2, a window of 1,024 residues centered on the mutation was used; for Prot-T5, a window of 512 residues centered on the mutation was used, reflecting a more conservative practical limit for that architecture. This windowing ensured that the mutated position was always contained within the input window, so that no variant was discarded due to its position in the sequence. The same windowing procedure was applied consistently across both the log-ratio and delta-embedding representations, for each model.

\subsection{Log-Ratio Representation}

For a mutation at position $i$, the conditional probabilities assigned by the model to the wild-type residue $a_{wt}$ and the mutant residue $a_{mut}$ were extracted. The mutation score was computed as:

\begin{equation}
    log\_Ratio = logP(a_{mut} | context) - logP(a_{wt} | context)
\end{equation}

This representation estimates how compatible the mutant residue is with the contextual sequence constraints learned by the model.

The two protein language models were not queried in an equivalent manner for this representation. For ESM-2, the wild-type sequence was passed to the model without masking, and conditional probabilities were obtained directly from the output logits at the mutated position (a pseudo-log-likelihood scoring approach). For Prot-T5, the residue at the mutated position was replaced with a sentinel token (\texttt{<extra\_id\_0>}), and the corresponding conditional probability was obtained from the decoder output (a masked marginal probability). Because the wild-type residue is visible to ESM-2 at scoring time but is masked from Prot-T5, log-ratio scores from the two models are not directly comparable in magnitude and were evaluated independently rather than benchmarked against one another in absolute terms.

\subsection{Delta-Embedding Representation}

Hidden-state embeddings corresponding to the mutated residue position were extracted for both the wild-type and mutant sequences, each passed through the model without masking. A delta-embedding vector was computed as:

\begin{equation}
    \Delta\epsilon = \epsilon_{mut} - \epsilon_{wt}
\end{equation}

where $\epsilon_{wt}$ and $\epsilon_{mut}$ denote the latent embeddings of the wild-type and mutant sequences at the mutated position, respectively. The zero-shot score used for evaluation was the L2 norm of this vector, $\lVert \Delta\epsilon \rVert_2$, under the rationale that a mutation producing a larger shift in the model's learned representation space corresponds to a stronger disruption of the sequence constraints captured during pretraining.

Embeddings were extracted from the final layer of each model's relevant module: layer 33 for ESM-2, accessed via the encoder-only \texttt{EsmModel} class, and layer 24 (final encoder layer) for Prot-T5, accessed via the \texttt{T5EncoderModel} class. The choice of the final layer for both models followed a criterion of methodological consistency rather than a search for an optimal layer; intermediate layers were not evaluated in this study and are noted as a direction for future work.

For Prot-T5 specifically, the delta-embedding representation was extracted differently from the log-ratio representation described above: rather than using the decoder with a masked sentinel token, the encoder-only model was used with the mutant residue present, unmasked, directly in the input sequence. As with the ESM-2/Prot-T5 log-ratio asymmetry noted above, the log-ratio and delta-embedding scores computed for Prot-T5 are therefore not extracted under comparable conditions and should not be compared in absolute magnitude; only ranking-based metrics (ROC-AUC, PR-AUC, Spearman correlation) are appropriate across these two representations.

Token-to-residue index alignment differs between the two models due to differences in special tokens: for ESM-2, the \texttt{<cls>} token occupies index 0, so the one-based protein position $i$ maps to index $i$ in the hidden-state tensor; for Prot-T5, no special token precedes the sequence (only an end-of-sequence token follows it), so position $i$ maps to index $i - 1$. Both index mappings were validated on a sample of variants by confirming that the residue retrieved at the computed embedding-extraction index matched the annotated wild-type amino acid.

\subsection{Dataset Construction}

The dataset was organized at the gene level, with each sample corresponding to a single human missense variant represented by the tuple \textbf{(gene, protein position, wild-type amino acid, mutant amino acid)}, together with a binary pathogenicity label. As described above, full protein reference sequences were retrieved from the Ensembl REST API \cite{ensembl} using transcript identifiers, rather than from UniProt.

Genes were stratified according to the number of annotated missense variants with both pathogenicity classes represented, in order to define an operational criterion for ``low-data'' evaluability. A gene was considered evaluable if it had at least five variants spanning both classes (2,253 genes); genes with three or four such variants were considered marginal (446 genes); genes with fewer than three variants, or with variants from a single class only, were excluded from per-gene evaluation (5,598 genes). Within the evaluable set, the median number of variants per gene was 3, the 75th percentile was 8, and the maximum was 859; this set comprised 54,848 variants across 3,130 unique proteins.

Processed datasets were serialized in Apache Parquet format to ensure efficient storage and reproducibility, while intermediate pre-processing and data integration steps were implemented using the Python Pandas ecosystem.

\subsection{Experimental Settings}

This study focuses on low-data regimes commonly observed in genes associated with rare diseases, operationalized through the gene-level evaluability criterion described above. Experiments were designed exclusively around the ability of pre-trained protein language models to infer mutation effects without any supervised, task-specific training: pathogenicity-related functional impact was inferred directly from pLM-derived mutation scores --- log-ratio and delta-embedding (L2 norm) --- with no downstream classifiers, fine-tuning, or use of population frequency as a model feature. The intrinsic representations learned by the pLMs were evaluated directly against ClinVar pathogenicity annotations on a per-gene basis.

\subsection{Evaluation Metrics}

For each of the four representation--model combinations (ESM-2 log-ratio, ESM-2 delta-norm, Prot-T5 log-ratio, Prot-T5 delta-norm), scores were evaluated per gene against the binary ClinVar label using ROC-AUC and PR-AUC, restricted to the evaluable gene set defined above. Pairwise comparisons between representations were assessed using Wilcoxon signed-rank tests on per-gene ROC-AUC values, paired by gene.

\subsection{Computational Environment}

All pre-processing and inference were implemented in Python using Pandas for data manipulation, NumPy for numerical operations, the Hugging Face Transformers library for pLM access, and SciPy/scikit-learn for evaluation and statistical testing. Variant annotation was performed with snpEff in a Bioconda-managed computational environment. All protein language model inference was performed on a MacBook Pro with an Apple M4 Max chip and 36 GiB of unified memory, using the PyTorch MPS backend (CUDA was not available); no fine-tuning was performed, and all reported representations correspond to inference-only computation.

\bibliographystyle{sbc}
\bibliography{references}

\end{document}
